% HAFF: Holographic Alaya-Field Framework
% General abstract: world, physical interface and subject formation.
\documentclass[12pt,a4paper]{article}
\usepackage{amsmath,amssymb}
\usepackage[hidelinks]{hyperref}
\usepackage{geometry}
\geometry{margin=1in}
\hypersetup{pdftitle={The Holographic Alaya-Field Framework: World, Interface and Subject},pdfauthor={Sidong Liu}}
\title{The Holographic Alaya--Field Framework:\\
World, Interface and Subject}
\author{Sidong Liu, PhD\\iBioStratix Ltd\\\texttt{sidongliu@hotmail.com}}
\date{October 2026}
\begin{document}
\emergencystretch=2em
\raggedbottom
\maketitle

\section*{General Abstract}

How do physical interactions form a stable world, the interfaces through which it can be reached, and systems whose histories organize their own participation in it? The Holographic Alaya--Field Framework studies these three questions together. Its central proposal is that usable structure depends on relations of access that can form, persist, be shared and be revised. A channel's resolution, a protected algebra, a finite task capability and a spacetime geometry are different objects. The theory connects them through specified interactions, records and reconstruction maps.

The world line begins with physical selection rather than an observer's arbitrary redescription. HAFF~A and B show how a supplied interaction distinguishes persistent blocks. Their coherence-decay rates define a Euclidean rate relation determined by the supplied coupling spectra, not a spatial geometry or a law of transfer between sectors. A response distance can be defined if separate experiments identify a reversible transfer generator; deriving that generator and its spatial interpretation remains part of the locality conjecture. The selected centre is not a tensor factorization. Channel classifications, finite probes and nonisometric correlation geometries support this programme by showing which distinctions an algebra alone misses; they are not substitutes for the origin of space \cite{HAFFA,HAFFB}.

HAFF~B and E supply dynamics and instruments to this account. Couplings can select stable blocks that carry both classical labels and protected quantum content. Their finite lifetime matters as much as the exact fixed algebra: a task's risk can vary continuously even when that algebra changes discontinuously. A probe's allowed effects, the records it actually acquires, and the effects that persist need not coincide. Environmental fragments can disclose block labels without disclosing the quantum content within them. The models give concrete conditions for repeatable and public reading \cite{HAFFB,HAFFE}.

A system can use records to maintain its ability to intervene. HAFF~C gives this capacity a body, history and resource budget: the same present report can warrant different actions, and learned diagnostic relations change later capability \cite{HAFFC}. HAFF~F distinguishes an evolution parameter, a recoverable event order and a directed record history. In an explicit reversible protocol, endpoint records identify an initially unknown permutation of exchange events; ordinary pointer copying alone can leave the complete endpoint state independent of that order. Acquiring records and renewing degraded correlations then have different entropy accounts. These results recover relations of before and after under specified dynamics and resource conditions, without pretending that a given evolution parameter has been generated \cite{HAFFF}.

HAFF~D asks how selected access and actual matter can constrain one geometric response. Its main route uses the same conserved source in a physical preparation and a modular consistency test. With the stated ball-family, frame, area-entropy and background conditions, those constraints give the linearized Einstein equation, rather than merely fitting an unrelated metric to data. C-II and Foundation supply conditional source-preparation and identification interfaces. The entropy--area matching and general spatial interpretation remain substantive inputs. A second route realizes source, response and later sensing in a finite semiclassical mean-action scalar/JT process. Full quantum material and clock constructions supply separate physical evidence; sustained fixed-capacity thermal updating and joint sharp heat export with effective return remain stronger conjectures. A universal gravitational mechanism has not been obtained by joining these different constructions in words \cite{HAFFD,Foundation,Consciousness}.

Q-RAIF gives the world--control relation its own six-paper argument. A derives Clifford and Lorentzian readout under local propagation and single-quadratic-shell conditions; a common cone needs further dynamical conditions. B connects centred field-pulse cost to the actual control generators: source-state-independent cost imposes a scalar-square condition, while identification with the propagation form additionally needs generator alignment. C places the representation question first. Discarding a continuous-spectrum control port restricts the exact multiplicative domain; retaining control within the closed propagation gives an exact representation and a conditional Lorentzian extension. These are two model classes, not an exhaustive classification of nature. Mass--Energy fixes the positive spectral branch and measurement window. Null Cluster tests null access within a matched source task. No View separates readout blindness from physical source selection and constrains local filters by stress conservation. These steps keep positive control cost, information geometry and Lorentzian geometry distinct while establishing conditional bridges between them \cite{QRAIF}.

T-DOME asks how history forms and revises a system's own participation. I gives task-specific memory requirements and an exact loss of usable gross feedback work when acquired history is compressed. II supplies a finite formation theorem: training on an unknown task--channel coupling selects a two-bit reading frame whose risk is below that of every preselected direction, even when the latter is given its exact regression coefficient. Both use one test read; retained memory and total cost are not matched, and the result does not distinguish a self-frame from general learned parameters. The apparatus and finite reading menu remain inputs. III asks what a formed self-frame must retain when its own access changes: diagnostic evidence, reliability and maintenance value can require different summaries. The memory--self--introspection route combines task-specific necessity, finite formation and learned revision, not a universal inevitability theorem \cite{TDOME}.

Foundation tests existence, sharing and return. I constructs stable hulls from a source, noise and readable seed, with conditions for an internal standard Cut. II distinguishes shared carrier, shared state, agreement and correctness, and supplies an authorized weak-mixture source test in its conformal instances. III gives a finite semiclassical mean-action scalar/JT return: an acquired record changes the apparatus source, and the response enters a later X-basis reading in the same history. The retained references, initial mass and geometric constraints are matched. The signal is of order $10^{-9}$, an existence result rather than a practical gain. Separate quantum-material renewal constructions retain different prepared banks and endpoints; they are not silently combined into this mean-action realization \cite{Foundation,QRAIF}.

Consciousness carries the first-person question. I relates internal standard Cuts to geometric data in the stated net class and distinguishes actual restrictions of access from reversible recoding. II uses physically prepared sources to identify and constrain higher-dimensional response under its field assumptions. III connects history-formed reading frames to causal self-diagnosis and proposes one two-layer bridge: geometric differences must reach the body's instruments, decoded records and possible actions before they can organize a perspective; a further empirical hypothesis relates this organization to reflective experience. Independent markers, sham interventions and competing report-equivalent mechanisms distinguish that proposal from simply naming functional improvement consciousness \cite{Consciousness}.

Imprint and Brain Anchor make history physically consequential. Carrying a trace, reading it through a reference, combining traces and using them in a later inquiry are different achievements. Imprint's scale models quantify source legibility and the value of retained history; joint relational content need not be recoverable from its separate marginals. The chain composes records, maintenance and action in specified common processes \cite{Imprint}. Brain Anchor supplies a body mechanism: sensorimotor history learns membership and orientation, changes routing and can install protection for held content. Learning and migration consume that content's lifetime, so better identification can mean worse survival. Its physical models and biological intervention hypotheses are distinguished from evidence for a neural quantum code \cite{Brain}.

HAFF~G returns to internal execution and the limits of unification. History-selected programs can alter protection and matter sources without every program becoming a central record. An invariant permitted algebra can make states indistinguishable under every finite adaptive record experiment, yet enlarging the permissions can break that equivalence. Such an operational boundary is not a universal final boundary of reality. The supplied executor and the formation of useful permissions must be distinguished when assessing explanatory completeness \cite{HAFFG}.

The programme therefore retains five coupled generative targets: a physical spatial dictionary, the relation between records and time order, world--control alignment, a common gravitational mechanism, and historically formed first-person organization. The papers provide new transfers and conditional mechanisms for these targets, while the Imprint Chain states their types, additional assumptions and common-realization status. Finite operational return does not imply a universe generated without inputs. The postscript returns to cosmic origin and closure through a conjecture about formation and renewal of relational structure, distinct from its conditional bounds on a body's participation. The central prospect remains that interactions form usable distinctions, histories reorganize access, and this organization can become a system's own revisable way of living in a world \cite{Imprint,HAFFPS}.

\bigskip
\noindent\textbf{Keywords}: physical access, quantum channels, stable structure, records, self-frames, maintained agency, shared worlds, physical feedback

\begin{thebibliography}{99}
\small
\item[] All numeric record identifiers below denote DOIs with the prefix \texttt{10.5281/zenodo.}; they identify the series records, not public uploads of these working revisions.
\bibitem{HAFFA}
S.~Liu, \emph{Physical Coarse-Graining and the Problem of Emergent Geometry: The Holographic Alaya-Field Framework}, (2026), series record DOI: 10.5281/zenodo.18361706.
\bibitem{HAFFB}
S.~Liu, \emph{Coupling, Records, and Stable Relations: Formation, Protected Content and the Lifetime of a Geometry}, (2026), series record DOI: 10.5281/zenodo.18367060.
\bibitem{HAFFC}
S.~Liu, \emph{Causation, Agency, and Existence: Historical Control and Maintained Intervention}, (2026), series record DOI: 10.5281/zenodo.18374805.
\bibitem{HAFFD}
S.~Liu, \emph{Selection, Maintenance and Gravitational Response: Modular Constraints and a Physical Return}, (2026), series record DOI: 10.5281/zenodo.18388881.
\bibitem{HAFFE}
S.~Liu, \emph{Measurement as Accessibility: Measurement Permissions, Stable Records and Public Readout}, (2026), series record DOI: 10.5281/zenodo.18400065.
\bibitem{HAFFF}
S.~Liu, \emph{Records, Order and Intervention: Conditions for a Temporal Direction}, (2026), series record DOI: 10.5281/zenodo.18400425.
\bibitem{HAFFG}
S.~Liu, \emph{Structural Limits of Unification: Internal Selection, Revisable Permissions and Physical Sources}, (2026), series record DOI: 10.5281/zenodo.18402907.
\bibitem{QRAIF}
S.~Liu, \emph{Q-RAIF}, six October 2026 working papers: A (18525876), B (18525890), C (18528934), Mass--Energy (20280951), Null Cluster (20283057), and No View (20284454).
\bibitem{TDOME}
S.~Liu, \emph{T-DOME}, October 2026 working trilogy: I, \emph{Memory as a Usable Resource} (DOI: 10.5281/zenodo.18574342); II, \emph{Reference Frames Selected by History} (18579703); III, \emph{Keeping a Self-Model Calibrated} (18591771).
\bibitem{Foundation}
S.~Liu, \emph{Foundation}, October 2026 working trilogy: I, \emph{From Source to Manifestation} (DOI: 10.5281/zenodo.19627786); II, \emph{Shared Manifestation} (19641372); III, \emph{The Perfuming Loop} (19642189).
\bibitem{Consciousness}
S.~Liu, \emph{Consciousness}, October 2026 working trilogy: I, \emph{Geometry Relative to a Cut} (DOI: 10.5281/zenodo.19607387); II, \emph{Higher-Dimensional Consistency of Geometry and Modular Flow} (19607543); III, \emph{Finite Self-Models and the Explanatory Gap} (19607724).
\bibitem{Imprint}
S.~Liu, \emph{Imprint}, four October 2026 working papers: \emph{The World as Imprint} (DOI: 10.5281/zenodo.20122718); \emph{Renormalization and Imprint Filtering} (20151451); \emph{Stable Hulls as Imprint Integration} (20151508); \emph{The Imprint Chain} (20151594).
\bibitem{Brain}
S.~Liu, \emph{The Brain as a Plastic Observable-Algebra Network: Learning Shared Access from Sensorimotor History}, (2026), series record DOI: 10.5281/zenodo.20044329.
\bibitem{HAFFPS}
S.~Liu, \emph{On the Closure of Structure: Origins, Persistence and Renewed Participation}, (2026), series record DOI: 10.5281/zenodo.18407367.
\end{thebibliography}
\end{document}
